% Original source: Zhou, physical PDF pp. 24--28 (printed pp. 8--12).
% Figures: Figure 2.1 is recreated with native TikZ; Figure 2.2 remains a native picture.
% The source's caret is retained as an infix exponentiation operator in the problem statement.

% ZHOU-SOURCE-PAGE: 24 PRINTED: 8
2, 3, 4 are lighter than 5, 6, 7, 8, as expressed using an ``<'' sign. There is no need to
explain the scenario that 1, 2, 3, 4 are heavier than 5, 6, 7, 8. (Why?\footnotemark[4])

If the two groups balance, this immediately tells us that the defective ball is in 9, 10, 11
and 12, and it is either lighter (\(L\)) or heavier (\(H\)) than other balls. Then we take 9,
10 and 11 from group 3 and compare balls 9, 10 with 8, 11. Here we have already figured out
that 8 is a normal ball. If 9, 10 are lighter, it must mean either 9 or 10 is \(L\) or 11 is
\(H\). In which case, we just compare 9 with 10. If 9 is lighter, 9 is the defective one and
it is \(L\); if 9 and 10 balance, then 11 must be defective and \(H\); If 9 is heavier, 10 is
the defective one and it is \(L\). If 9, 10 and 8, 11 balance, 12 is the defective one. If 9,
10 is heavier, than either 9 or 10 is \(H\), or 11 is \(L\).

You can easily follow the tree in Figure 2.1 for further analysis and it is clear from the
tree that all possible scenarios can be resolved in 3 measurements.

\begin{figure}[t]
\centering
\begin{tikzpicture}[
  x=1mm,
  y=1mm,
  >=Stealth,
  font=\scriptsize,
  line cap=round,
  line join=round,
  pan/.style={draw, ellipse, minimum height=5.5mm, inner xsep=2.5mm, inner ysep=0.7mm},
  branch/.style={font=\scriptsize, align=center, inner sep=0.5pt, fill=white},
  result/.style={font=\scriptsize, inner sep=0.2pt}
]
\path[use as bounding box] (0,0) rectangle (151,89);

% First weighing and its two nontrivial branches.
\node[pan, minimum width=7mm] (root) at (75,86) {?};
\coordinate (rootpivot) at (75,80);
\coordinate (rootfork) at (75,68);
\draw[->] (root.south) -- (rootpivot);
\node at (72,82.6) {$+$};
\node at (78,80.8) {$-$};
\node[pan, minimum width=24mm] (firstleft) at (52,74) {1,2,3,4};
\node[pan, minimum width=24mm] (firstright) at (99,74) {5,6,7,8};
\draw (rootpivot) -- (firstleft.north east);
\draw (rootpivot) -- (firstright.north west);
\draw (rootpivot) -- (rootfork);

\coordinate (leftentry) at (35,58);
\coordinate (rightentry) at (116,58);
\draw[->] (rootfork) -- (leftentry);
\draw[->] (rootfork) -- (rightentry);
\node[branch, text width=47mm] at (43,64) {1/2/3/4 $L$ or 5/6/7/8 $H$};
\node[branch, text width=44mm] at (108,64) {9/10/11/12 $L$ or $H$};
\node at (58,69.5) {$-$};
\node at (92,69.5) {$=$};

% Left branch: compare 1,2,5 with 3,6,9.
\node[pan, minimum width=19mm] (leftpanone) at (22,52) {1,2,5};
\node[pan, minimum width=19mm] (leftpantwo) at (48,52) {3,6,9};
\draw (leftentry) -- (leftpanone.north east);
\draw (leftentry) -- (leftpantwo.north west);
\coordinate (leftfork) at (35,44);
\draw (leftfork) -- (leftpanone.south east);
\draw (leftfork) -- (leftpantwo.south west);

\node[pan, minimum width=16mm] (leftcompareone) at (9,27) {1 $-$ 2};
\node[pan, minimum width=16mm] (leftcomparetwo) at (35,27) {7 $-$ 8};
\node[pan, minimum width=16mm] (leftcomparethree) at (61,27) {3 $-$ 9};
\draw[->] (leftfork) -- (leftcompareone.north);
\draw[->] (leftfork) -- (leftcomparetwo.north);
\draw[->] (leftfork) -- (leftcomparethree.north);
\node[branch, text width=21mm] at (9,39.5) {$1/2L$ or $6H$};
\node[branch, text width=21mm] at (35,39.5) {$4L$ or $7/8H$};
\node[branch, text width=18mm] at (61,39.5) {$5H$ or $3L$};
\node at (9,42.5) {$-$};
\node at (35,42.5) {$=$};
\node at (61,42.5) {$+$};

\draw[->] (leftcompareone.south) -- (3,10);
\draw[->] (leftcompareone.south) -- (9,10);
\draw[->] (leftcompareone.south) -- (15,10);
\node at (6,19) {$-$};
\node at (9,18) {$=$};
\node at (12,19) {$+$};
\node[result] at (3,7) {$1L$};
\node[result] at (9,7) {$6H$};
\node[result] at (15,7) {$2L$};

\draw[->] (leftcomparetwo.south) -- (29,10);
\draw[->] (leftcomparetwo.south) -- (35,10);
\draw[->] (leftcomparetwo.south) -- (41,10);
\node at (32,19) {$-$};
\node at (35,18) {$=$};
\node at (38,19) {$+$};
\node[result] at (29,7) {$8H$};
\node[result] at (35,7) {$4L$};
\node[result] at (41,7) {$7H$};

\draw[->] (leftcomparethree.south) -- (55,10);
\draw[->] (leftcomparethree.south) -- (61,10);
\draw[->] (leftcomparethree.south) -- (67,10);
\node at (58,19) {$-$};
\node at (61,18) {$=$};
\node at (64,19) {$+$};
\node[result] at (55,7) {$3L$};
\node[result] at (61,7) {$=$};
\node[result] at (67,7) {$5H$};

% Right branch: compare 9,10 with 8,11.
\node[pan, minimum width=19mm] (rightpanone) at (103,52) {9,10};
\node[pan, minimum width=19mm] (rightpantwo) at (129,52) {8,11};
\draw (rightentry) -- (rightpanone.north east);
\draw (rightentry) -- (rightpantwo.north west);
\coordinate (rightfork) at (116,44);
\draw (rightfork) -- (rightpanone.south east);
\draw (rightfork) -- (rightpantwo.south west);

\node[pan, minimum width=16mm] (rightcompareone) at (88,27) {9 $-$ 10};
\node[pan, minimum width=16mm] (rightcomparetwo) at (115,27) {8 $-$ 12};
\node[pan, minimum width=16mm] (rightcomparethree) at (142,27) {9 $-$ 10};
\draw[->] (rightfork) -- (rightcompareone.north);
\draw[->] (rightfork) -- (rightcomparetwo.north);
\draw[->] (rightfork) -- (rightcomparethree.north);
\node[branch, text width=21mm] at (88,39.5) {$9/10L$ or $11H$};
\node[branch, text width=21mm] at (115,39.5) {$12L$ or $12H$};
\node[branch, text width=21mm] at (142,39.5) {$9/10H$ or $11L$};
\node at (88,42.5) {$-$};
\node at (115,42.5) {$=$};
\node at (142,42.5) {$+$};

\draw[->] (rightcompareone.south) -- (82,10);
\draw[->] (rightcompareone.south) -- (88,10);
\draw[->] (rightcompareone.south) -- (94,10);
\node at (85,19) {$-$};
\node at (88,18) {$=$};
\node at (91,19) {$+$};
\node[result] at (82,7) {$9L$};
\node[result] at (88,7) {$11H$};
\node[result] at (94,7) {$10L$};

\draw[->] (rightcomparetwo.south) -- (109,10);
\draw[->] (rightcomparetwo.south) -- (121,10);
\node at (112,19) {$-$};
\node at (118,19) {$+$};
\node[result] at (109,7) {$12H$};
\node[result] at (121,7) {$12L$};

\draw[->] (rightcomparethree.south) -- (136,10);
\draw[->] (rightcomparethree.south) -- (142,10);
\draw[->] (rightcomparethree.south) -- (148,10);
\node at (139,19) {$-$};
\node at (142,18) {$=$};
\node at (145,19) {$+$};
\node[result] at (136,7) {$10H$};
\node[result] at (142,7) {$11L$};
\node[result] at (148,7) {$9H$};
\end{tikzpicture}
\caption{Tree diagram to identify the defective ball in 12 balls}
\label{fig:zhou-figure-2-1}
\end{figure}

\footnotetext[4]{Here is where the symmetry idea comes in. Nothing makes the 1, 2, 3, 4 or
5, 6, 7, 8 labels special. If 1, 2, 3, 4 are heavier than 5, 6, 7, 8, let's just exchange the
labels of these two groups. Again we have the case of 1, 2, 3, 4 being lighter than 5, 6, 7, 8.}

In general if you have the information as to whether the defective ball is heavier or
% ZHOU-SOURCE-PAGE: 25 PRINTED: 9
lighter, you can identify the defective ball among up to \(3^n\) balls using no more than \(n\)
measurements since each weighing reduces the problem size by 2/3. If you have no information
as to whether the defective ball is heavier or lighter, you can identify the defective ball among
up to \((3^n-3)/2\) balls using no more than \(n\) measurements.

\end{solution}

\begin{problem}{Trailing zeros}
How many trailing zeros are there in 100! (factorial of 100)?
\end{problem}

\begin{solution}
This is an easy problem. We know that each pair of 2 and 5 will give a trailing zero. If we
perform prime number decomposition on all the numbers in 100!, it is obvious that the frequency
of 2 will far outnumber the frequency of 5. So the frequency of 5 determines the number of
trailing zeros. Among numbers 1,2,\ldots,99, and 100, 20 numbers are divisible by 5 (5, 10,
\ldots, 100). Among these 20 numbers, 4 are divisible by \(5^2\) (25, 50, 75, 100). So the
total frequency of 5 is 24 and there are 24 trailing zeros.
\end{solution}

\begin{problem}{Horse race}
There are 25 horses, each of which runs at a constant speed that is different from the other
horses\textquotesingle{}. Since the track only has 5 lanes, each race can have at most 5 horses.
If you need to find the 3 fastest horses, what is the minimum number of races needed to identify
them?
\end{problem}

\begin{solution}
This problem tests your basic analytical skills. To find the 3 fastest horses, surely all horses
need to be tested. So a natural first step is to divide the horses to 5 groups (with horses 1-5,
6-10, 11-15, 16-20, 21-25 in each group). After 5 races, we will have the order within each
group, let's assume the order follows the order of numbers (e.g., 6 is the fastest and 10 is the
slowest in the 6-10 group)\footnotemark[5]. That means 1, 6, 11, 16 and 21 are the fastest
within each group.

Surely the last two horses within each group are eliminated. What else can we infer? We know
that within each group, if the fastest horse ranks 5th or 4th among 25 horses, then all horses in
that group cannot be in top 3; if it ranks the 3rd, no other horse in that group can be in the top
3; if it ranks the 2nd, then one other horse in that group may be in top 3, if it ranks the first,
then two other horses in that group may be in top 3.

\footnotetext[5]{Such an assumption does not affect the generality of the solution. If the order
is not as described, just change the labels of the horses.}

% ZHOU-SOURCE-PAGE: 26 PRINTED: 10
So let's race horses 1, 6, 11, 16 and 21. Again without loss of generality, let's assume the
order is 1, 6, 11, 16 and 21. Then we immediately know that horses 4-5, 8-10, 12-15, 16-20
and 21-25 are eliminated. Since 1 is fastest among all the horses, 1 is in. We need to determine
which two among horses 2, 3, 6, 7 and 11 are in top 3, which only takes one extra race.

So all together we need 7 races (in 3 rounds) to identify the 3 fastest horses.
\end{solution}

\begin{problem}{Infinite sequence}
If \(x\mathbin{\text{\textasciicircum}}x\mathbin{\text{\textasciicircum}}x
\mathbin{\text{\textasciicircum}}x\mathbin{\text{\textasciicircum}}x
\mathbin{\text{\textasciicircum}}\cdots=2\), where
\(x\mathbin{\text{\textasciicircum}}y=x^y\), what is \(x\)?
\end{problem}

\begin{solution}
This problem appears to be difficult, but a simple analysis will give an elegant solution. What
do we have from the original equation?

\[
\begin{aligned}
\lim_{n\to\infty}
  \underbrace{x^{\left(x^{\left(x^{\cdots^x}\right)}\right)}}_{n\ \mathrm{terms}}
  &=2\\
  &\Longleftrightarrow\quad
  \lim_{n\to\infty}
  \underbrace{x^{\left(x^{\left(x^{\cdots^x}\right)}\right)}}_{n-1\ \mathrm{terms}}
  =2.
\end{aligned}
\]

In other words, as \(n\to\infty\), adding or minus one
\(x\mathbin{\text{\textasciicircum}}\) should yield the same result.

So
\[
x^{x^{x^{x^{\cdots}}}}=x^{\left(x^{x^{x^{\cdots}}}\right)}=x^2=2
\quad\Longrightarrow\quad x=\sqrt{2}.
\]
\end{solution}

\subsection{Thinking Out of the Box}

\begin{problem}{Box packing}
Can you pack 53 bricks of dimensions \(1\times1\times4\) into a \(6\times6\times6\) box?
\end{problem}

\begin{solution}
This is a nice problem extended from a popular chess board problem. In that problem, you have a
\(8\times8\) chess board with two small squares at the opposite diagonal corners removed. You
have many bricks with dimension \(1\times2\). Can you pack 31 bricks into the remaining 62 squares?
(An alternative question is whether you can cover all 62 squares using bricks without any bricks
overlapping with each other or sticking out of the board, which requires a similar analysis.)

A real chess board figure surely helps the visualization. As shown in Figure 2.2, when a chess
board is filled with alternative black and white squares, both squares at the opposite diagonal
corners have the same color. If you put a $1\times2$ brick on the board, it will always cover one
black square and one white square. Let's say it's the two black corner squares were removed, then
the rest of the board can fit at most 30 bricks since we only have 30 black squares left (and each
brick requires one black square). So to pack 31 bricks is out of the question. To cover all 62
squares without overlapping or overreaching, we must have exactly 31 bricks. Yet we have proved
that 31 bricks cannot

% ZHOU-SOURCE-PAGE: 27 PRINTED: 11
fit in the 62 squares left, so you cannot find a way to fill in all 62 squares without overlapping
or overreaching.

\begin{figure}[t]
\centering
{\setlength{\unitlength}{5mm}
\begin{picture}(17,13)
  \put(5,3){\framebox(8,8){}}
  \multiput(7,3)(2,0){3}{\rule{\unitlength}{\unitlength}}
  \multiput(6,4)(2,0){4}{\rule{\unitlength}{\unitlength}}
  \multiput(5,5)(2,0){4}{\rule{\unitlength}{\unitlength}}
  \multiput(6,6)(2,0){4}{\rule{\unitlength}{\unitlength}}
  \multiput(5,7)(2,0){4}{\rule{\unitlength}{\unitlength}}
  \multiput(6,8)(2,0){4}{\rule{\unitlength}{\unitlength}}
  \multiput(5,9)(2,0){4}{\rule{\unitlength}{\unitlength}}
  \multiput(6,10)(2,0){3}{\rule{\unitlength}{\unitlength}}
  \put(5.5,3.5){\circle{2}}
  \put(12.5,10.5){\circle{2}}
  \put(14,10.5){\vector(1,0){2}}
  \put(16.5,10.5){\makebox(0,0)[l]{Removed}}
  \put(4,3.5){\vector(-1,0){2}}
  \put(2,3.5){\makebox(0,0)[r]{Removed}}
\end{picture}}
\caption{Chess board with alternative black and white squares}
\label{fig:zhou-figure-2-2}
\end{figure}

Just as any good trading strategy, if more and more people get to know it and replicate it, the
effectiveness of such a strategy will disappear. As the chess board problem becomes popular,
many interviewees simply commit it to memory (after all, it's easy to remember the answer). So
some ingenious interviewer came up with the newer version to test your thinking process, or at
least your ability to extend your knowledge to new problems.

If we look at the total volume in this 3D problem, 53 bricks have a volume of 212, which is
smaller then the box's volume 216. Yet we can show it is impossible to pack all the bricks into
the box using a similar approach as the chess board problem. Let's imagine that the \(6\times6\times6\)
box is actually comprised of small \(2\times2\times2\) cubes. There should be 27 small cubes.
Similar to the chess board (but in 3D), imagine that we have black cubes and white cubes
alternates---it does take a little 3D visualization. So we have either 14 black cubes \& 13 white
cubes or 13 black cubes \& 14 white cubes. For any \(1\times1\times4\) brick that we pack into the
box, half (\(1\times1\times2\)) of it must be in a black \(2\times2\times2\) cube and the other
half must be in a white \(2\times2\times2\) cube. The problem is that each \(2\times2\times2\)
cube can only be used by 4 of the \(1\times1\times4\) bricks. So for the color with 13 cubes,
be it black or white, we can only use them for 52 \(1\times1\times4\) tubes. There is no way to
place the 53th brick. So we cannot pack 53 bricks of dimensions \(1\times1\times4\) into a
\(6\times6\times6\) box.

\end{solution}

\begin{problem}{Calendar cubes}
You just had two dice custom-made. Instead of numbers 1 -- 6, you place single-digit numbers
on the faces of each dice so that every morning you can arrange the dice in a way as to make the
two front faces show the current day of the month. You must use both dice (in other words, days
1 -- 9 must be shown as 01 -- 09), but you can switch the

% ZHOU-SOURCE-PAGE: 28 PRINTED: 12
order of the dice if you want. What numbers do you have to put on the six faces of each of the two
dice to achieve that?
\end{problem}

\begin{solution}
The days of a month include 11 and 22, so both dice must have 1 and 2. To express single-digit
days, we need to have at least a 0 in one dice. Let's put a 0 in dice one first. Considering that
we need to express all single digit days and dice two cannot have all the digits from 1 -- 9, it's
necessary to have a 0 in dice two as well in order to express all single-digit days.

So far we have assigned the following numbers:

\begin{center}
{\renewcommand{\arraystretch}{1.15}%
\setlength{\tabcolsep}{5pt}%
\begin{tabular}{@{}lcccccc@{}}
\toprule
Dice one & 1 & 2 & 0 & ? & ? & ? \\
\midrule
Dice two & 1 & 2 & 0 & ? & ? & ? \\
\bottomrule
\end{tabular}}
\end{center}

If we can assign all the rest of digits 3, 4, 5, 6, 7, 8, and 9 to the rest of the faces, the
problem is solved. But there are 7 digits left. What can we do? Here's where you need to think
out of the box. We can use a 6 as a 9 since they will never be needed at the same time! So,
simply put 3, 4, and 5 on one dice and 6, 7, and 8 on the other dice, and the final numbers on
the two dice are:

\begin{center}
{\renewcommand{\arraystretch}{1.15}%
\setlength{\tabcolsep}{5pt}%
\begin{tabular}{@{}lcccccc@{}}
\toprule
Dice one & 1 & 2 & 0 & 3 & 4 & 5 \\
\midrule
Dice two & 1 & 2 & 0 & 6 & 7 & 8 \\
\bottomrule
\end{tabular}}
\end{center}
\end{solution}

\begin{problem}{Door to offer}
You are facing two doors. One leads to your job offer and the other leads to exit. In front of
either door is a guard. One guard always tells lies and the other always tells the truth. You can
only ask one guard one yes/no question. Assuming you do want to get the job offer, what question
will you ask?
\end{problem}

\begin{solution}
This is another classic brain teaser (maybe a little out-of-date in my opinion). One popular
answer is to ask one guard: ``Would the other guard say that you are guarding the door to the
offer?'' If he answers yes, choose the other door; if he answers no, choose the door this guard is
standing in front of.

There are two possible scenarios:

\begin{enumerate}
\item Truth teller guards the door to offer; Liar guards the door to exit.
\item Truth teller guards the door to exit; Liar guards the door to offer.
\end{enumerate}

If we ask a guard a direct question such as ``Are you guarding the door to the offer?'' For
scenario 1, both guards will answer yes; for scenario 2, both guards will answer no. So a
